% ===== §13 =====
\section{The Capacities, not a Procedure}
\label{p27:sec:capacities}

\noindent
The model this paper proposes is a set of capacities and not a procedure, and the distinction follows as a consequence of the account and is no matter of presentation. A procedure is an ordered sequence of steps, each following the last, the whole completed when the final step is reached. A model of relational understanding cast as a procedure would name a first thing to do, then a second, and so on to a last, and would represent understanding as achieved when the sequence is run. The account developed here forbids this, for the same reason it forbids the rehearsal. The steps of such a procedure would have to follow one another in time, the earlier completed before the later begins, and yet the matters they concern do not order themselves so. The value of a thing given may return only after the later steps have long since passed, the feedback that a procedure would require at each step to confirm its completion before the next is delayed and uncertain, and the whole may span a length of time across which no sequence could be held as a single connected run. A procedure presupposes that its steps are adjacent in time and that each completes before the next begins. The generativity of relational value, which may disclose the worth of an early act only in a late season of the bond, breaks this presupposition, and with it the possibility of a procedure.

What the account permits, in the place of a procedure, is a set of capacities. A capacity is a formed power that can be drawn upon at any moment the relation requires it, in any order, and whose exercise at one moment does not wait upon the completion of another. Capacities can be held together, brought to bear as the unrepeatable and slow-returning history of a bond calls for them, without the temporal adjacency a procedure demands. To cast the model as capacities and not as steps is thus to make it consistent with the framework's own conception of relational time, on which the acts of understanding are drawn, each in its moment, from powers the person has formed, and do not line up in a sequence, from powers the person has formed and holds in readiness.

The model names seven such capacities. They are not stages and are not ordered, and the numbering that follows is for reference and not for sequence. The first is the recognition of the other's language world, the acknowledgment that the other renders the world in a symbolic system with its own generative rules and not merely in a different vocabulary. The second is the entry into that world, the capacity to suspend one's own system far enough to let the other's renderings become natural, tested by whether one can generate a rendering the other affirms as their own. The third is the construction of a structural mapping, the observation of how the relations that hold in one's own world correspond, at the level of structure, to the relations that hold in the other's. The fourth is translation, the rendering of a matter from the other's world into one's own with the loss marked and the mapping honored. The fifth is the throwing of a generative thing into the other's own unfolding, the offering, in the other's terms, of something from which the other generates what they had not observed, as against the forcing upon the other of a conclusion of one's own. The sixth is a non-intervention, the capacity to withhold one's steering and let a shared dynamic develop between the two. The seventh is the joint recognition and preservation of what is generated, the noticing together of a generated matter's value, the finding of where the two worlds meet, and the keeping of a relational experience against the season in which its value returns.

\begin{claim}[The model: capacities in two orders of cultivation]
Relational understanding in intimacy is proposed as a set of seven capacities and not a procedure, since the generativity of relational value forbids the temporal adjacency a procedure requires. The seven are the recognition of the other's world, the entry into it, the construction of a structural mapping, translation, the throwing of a generative thing into the other's unfolding, non-intervention, and the joint recognition and preservation of what is generated. They divide into two orders. The first four are near-practicable, cultivable in the loop of attempt and correction in fields outside the bond, and brought to the bond as developed skills. The last three reach into the wager, give feedback only after delay and into an unrepeatable history, and cannot be honed by a loop; they are cultivated as a standing disposition and not as skills. Non-possession is the condition of all seven, without which they decline into a refined extraction.
\end{claim}

\begin{figure}[h]
\centering
\begin{tikzpicture}[
  every node/.style={font=\small},
  capnode/.style={draw=rose, line width=0.7pt, rounded corners=2pt, fill=white,
              inner sep=4pt, minimum height=8mm, text width=3.5cm, align=center, text=inksoft},
  ordlabel/.style={font=\footnotesize\itshape, text=gold, align=center}
]
% enclosing field: non-possession
\draw[gold, line width=0.9pt, rounded corners=10pt]
  (-6.6,-4.7) rectangle (6.6,4.2);
\node[text=gold, font=\footnotesize\bfseries] at (0,3.75)
  {the condition of non-possession, enclosing all seven};

% ---- near-practicable region (left column) ----
\node[ordlabel, text width=4.4cm] at (-3.4,3.0) {near-practicable\\cultivable in the loop};
\node[capnode] (c1) at (-3.4,1.9)  {1\ \ recognition of the other's world};
\node[capnode] (c2) at (-3.4,0.5)  {2\ \ entry into that world};
\node[capnode] (c3) at (-3.4,-0.9) {3\ \ structural mapping};
\node[capnode] (c4) at (-3.4,-2.3) {4\ \ translation};
\node[ordlabel, text width=4.4cm] at (-3.4,-3.65)
  {held in readiness and drawn on in any order, not in sequence};

% dividing line between the two orders
\draw[gold, dashed, line width=0.5pt] (0,-4.2) -- (0,3.4);

% ---- wagered region (right column) ----
\node[ordlabel, text width=4.4cm] at (3.4,3.0) {wagered\\cultivable only as a disposition};
\node[capnode] (c5) at (3.4,1.9)  {5\ \ throwing a generative thing};
\node[capnode] (c6) at (3.4,0.5)  {6\ \ non-intervention};
\node[capnode] (c7) at (3.4,-0.9) {7\ \ joint recognition and preservation};

% 5 and 6 as one gesture, tie on the right side (no directional arrow)
\draw[rose, line width=0.5pt] (c5.east) .. controls (5.9,1.2) .. (c6.east);
\node[ordlabel, text=rose, text width=1.4cm] at (6.05,1.2) {one gesture};
\end{tikzpicture}
\caption{The model as capacities, not a procedure. The seven capacities are held in readiness and drawn upon in any order, with no sequence among them. The four on the left admit cultivation in the loop of attempt and correction; the three on the right reach into the wager and are cultivated only as a standing disposition. The fifth and sixth compose a single gesture in two moments. Non-possession is drawn as the field enclosing all seven, the condition under which each is what it is, and not as an eighth capacity.}
\label{p27:fig:capacities}
\end{figure}

The first four capacities, the recognition, the entry, the mapping, and the translation, were treated at their root in the companion work and are near-practicable in the sense the previous part gave. They are the crossing of worlds in its cognitive labor, and while intimacy raises them to their highest difficulty, they can be developed in the wider life where the crossing is correctable, and brought to the bond as powers already formed. This paper does not add to their theory beyond what the domain of intimacy requires, and turns instead to the three that reach into the wager, for it is there that intimacy makes its distinctive demand and there that the account has most to add. The next sections take the fifth and the seventh in turn, the throwing of a generative thing and the preservation of relational experience, with the sixth, non-intervention, drawn out alongside the fifth, since the two are a single gesture in two moments. The condition of non-possession, named in the claim as the ground of all seven, is drawn together at the last.
